% ============================================================
%  preamble.tex — Wasserstein 距離セミナー用 共通設定
%  Villani (2009) を主たる文献とし、記法もこれに合わせる
% ============================================================

% --- 基本パッケージ ---
\usepackage{amsmath,amssymb,amsthm,mathtools}
\usepackage[dvipdfmx,hidelinks]{hyperref}
\usepackage{pxjahyper}

% --- 図・表 ---
\usepackage{tikz}
\usepackage{float}
\usetikzlibrary{decorations.pathreplacing,calc,arrows.meta,patterns,positioning}
\usepackage{graphicx}
\graphicspath{{fig/}}

% --- 定理環境 (tcolorbox) ---
\usepackage{tcolorbox}
\tcbuselibrary{theorems,breakable}

% Definition: 青系
\newtcbtheorem[number within=section]{definition}{Def.}{
  colback=blue!5, colframe=blue!40!black, fonttitle=\bfseries, breakable
}{def}
% Proposition: オレンジ系
\newtcbtheorem[use counter from=definition]{proposition}{Prop.}{
  colback=orange!5, colframe=orange!50!black, fonttitle=\bfseries, breakable
}{prop}
% Theorem: 赤系（章・節の目玉となる重要結果）
\newtcbtheorem[use counter from=definition]{theorem}{Thm.}{
  colback=red!5, colframe=red!40!black, fonttitle=\bfseries, breakable
}{thm}
% Claim: 緑系（Thm ほど深くはないが，証明を要する主張）
\newtcbtheorem[use counter from=definition]{claim}{Claim}{
  colback=green!5, colframe=green!45!black, fonttitle=\bfseries, breakable
}{clm}
% Lemma: シアン系（命題の証明に用いる補助結果）
\newtcbtheorem[use counter from=definition]{lemma}{Lem.}{
  colback=cyan!5, colframe=cyan!50!black, fonttitle=\bfseries, breakable
}{lem}
% Remark: 黄系
\newtcbtheorem[use counter from=definition]{remark}{Rem.}{
  colback=yellow!5, colframe=yellow!40!black, fonttitle=\bfseries, breakable
}{rem}
% Example: 灰色系
\newtcbtheorem[use counter from=definition]{example}{Ex.}{
  colback=gray!5, colframe=gray!50!black, fonttitle=\bfseries, breakable
}{ex}
% Corollary: 紫系
\newtcbtheorem[use counter from=definition]{corollary}{Cor.}{
  colback=purple!5, colframe=purple!40!black, fonttitle=\bfseries, breakable
}{cor}
% Memo: 薄い黒系（番号なし，発展的補足）
\newtcolorbox{memo*}{
  colback=black!7, colframe=black!40,
  fonttitle=\bfseries, title={Memo},
  breakable
}
% 証明環境の装飾
\tcolorboxenvironment{proof}{
  colback=white, colframe=black!20,
  leftrule=2pt, rightrule=0pt, toprule=0pt, bottomrule=0pt,
  breakable
}

% ブロックの機械可読メタデータ（proofgraph ツールが解析する）。
\newcommand{\blockmeta}[1]{}

% --- 便利マクロ ---
\newcommand{\R}{\mathbb{R}}
\newcommand{\N}{\mathbb{N}}
\newcommand{\abs}[1]{\lvert #1 \rvert}
\newcommand{\norm}[1]{\lVert #1 \rVert}

% --- 最適輸送の記法（Villani 準拠）---
\newcommand{\Borel}{\mathcal{B}}
\newcommand{\Prob}{\mathcal{P}}
\newcommand{\Pp}[1]{\mathcal{P}_{#1}}
\newcommand{\Couplings}{\Pi}
\newcommand{\Wass}{W}
\newcommand{\Id}{\mathrm{Id}}
\renewcommand{\d}{\mathrm{d}}               % 積分の d

% --- サイト用デモブロック指示（PDF では無効，tex2md が :::demo に変換する）---
\newcommand{\demohint}[1]{}
